AgentSense aims to provide a realistic social intelligence benchmark with enhanced diversity and complexity.
Following the Dramaturgical Theory, we propose an overall framework as in Figure~\ref{fig:fm}.

\paragraph{Scenarios}
The core component of AgentSense is the social scenario set, extracted from real-world scripts to guide and evaluate social interactions between agents.
A social scenario serves as a hypothetical context for simulating and analyzing social interactions, where two key components are measured:
(1) \textbf{Social Goal} is what the agent aims to achieve, such as resolving an issue or building a relationship.
The agent's proactive drive in social interactions, guided by this social goal, directs its \emph{active participation} in social dynamics.
(2) \textbf{Private Information} is information that is known solely to the agent and not to others.
The agent is tasked with inferring others' private information without explicitly inquiring about it, a process referred to as \emph{passive reasoning} during interactions.
In summary, an agent's social intelligence is reflected in its ability to pursue social goals while safeguarding private information, balancing active engagement with passive respect for individual privacy.

\paragraph{Scenario Templates}
Social scenarios in scripts always have a fixed group of characters, causing a lack of diversity.
To address this issue, we wipe out irrelevant character details to obtain \textbf{scenario templates}, which contains only background information and predefined character slots.
We can instantiate multiple scenarios from a scenario template by filling in the slots with different sets of synthesized characters satisfying the template's constraints.

\paragraph{Benchmarking}
After building scenarios from the extracted templates, benchmarking LLMs with AgentSense comes as follows:
(1) \textbf{Simulation}: We prompt the models to role-play the characters and interact with each other, trying to achieve their social goals.
(2) \textbf{Evaluation:} We evaluate the goal completion status of each model by interviewing the participants and third-party judges.
We also assess the model's implicit reasoning performance with multiple-choice questions.
